% GENERATED FILE. Edit canonical lessons, then run npm run generate:books.
% canonical-source-hash: fnv1a64:454977c4d332fbf5
% canonical-lessons: SA-C71-runs, SA-C71-walks, SA-C71-sits, SA-C71-standsup, SA-C71-falls

\chapter{Run, Walk, Sit, Stand}
\label{ch:sa-run-walk-sit-stand}
\hlchaptermodality{71}

\begin{chapteropening}
\textbf{By the end of this chapter:} \emph{I can say someone runs, walks, sits down, stands up or falls.}

Complete the last lesson of chapter 71: I can say someone runs, walks, sits down, stands up or falls.
\end{chapteropening}

\section[\texorpdfstring{\sk{धावति} (dhāvati)}{धावति (dhāvati)}]{\sk{धावति} (dhāvati) — runs}

\label{lesson:SA-C71-runs}



\begin{quote}
Before the new one: say the Sanskrit for the shoulder, then the Sanskrit for the lip.
\end{quote}

\subsection*{You'll want to know: \sk{धावति}}
\textbf{\sk{धावति}} (\emph{dhāvati}) — \textquotedblleft{}(he, she) runs\textquotedblright{}. It ends in \textbf{-\sk{ति}} like \textbf{\sk{गच्छति}}: one person, doing it now.

\begin{cousinweb}
\emph{Dhāvati} is the present of the root √dhāv, \textquotedblleft{}to run, flow\textquotedblright{}, already in the Rigveda (MW). CDIAL 6802 follows it into Hindi \emph{dhāwnā}, \textquotedblleft{}to run\textquotedblright{}.
\end{cousinweb}

\begin{grammarlens}[title={What you've built}]
A verb in the frame you know.
\end{grammarlens}

\subsection*{Your turn}
\begin{itemize}
\raggedright
  \item \emph{Say it:} \emph{dhāvati}
  \item \emph{Say it:} \emph{dhāvati}, once more
  \item \emph{Say it:} say \emph{bālaḥ dhāvati}
  \item \emph{Recall:} say the Sanskrit for the shoulder, then the Sanskrit for the lip, then say \emph{dhāvati} again
\end{itemize}

\begin{tcolorbox}[breakable,colback=teal!4,colframe=teal!35!black,title={Before you move on}]
What does \sk{धावति} mean? (\textquotedblleft{}Runs\textquotedblright{}.) Say it once more.
\end{tcolorbox}
\section[\texorpdfstring{\sk{चलति} (calati)}{चलति (calati)}]{\sk{चलति} (calati) — walks, moves}

\label{lesson:SA-C71-walks}



\begin{quote}
Before the new one: say the Sanskrit for the lip, then the Sanskrit for he runs.
\end{quote}

\subsection*{You'll want to know: \sk{चलति}}
\textbf{\sk{चलति}} (\emph{calati}) — \textquotedblleft{}(he, she) moves, walks\textquotedblright{}.

\begin{cousinweb}
\emph{Calati} is the present of the root √cal, \textquotedblleft{}to stir, move\textquotedblright{} (MW). It lives on as Hindi \emph{calnā}, \textquotedblleft{}to move, walk\textquotedblright{} (CDIAL 4715).
\end{cousinweb}

\begin{grammarlens}[title={What you've built}]
The slower way to go.
\end{grammarlens}

\subsection*{Your turn}
\begin{itemize}
\raggedright
  \item \emph{Say it:} \emph{calati}
  \item \emph{Say it:} \emph{calati}, once more
  \item \emph{Say it:} say \emph{saḥ calati}
  \item \emph{Recall:} say the Sanskrit for the lip, then the Sanskrit for he runs, then say \emph{calati} again
\end{itemize}

\begin{tcolorbox}[breakable,colback=teal!4,colframe=teal!35!black,title={Before you move on}]
What does \sk{चलति} mean? (\textquotedblleft{}Walks, moves\textquotedblright{}.) Say it once more.
\end{tcolorbox}
\section[\texorpdfstring{\sk{उपविशति} (upaviśati)}{उपविशति (upaviśati)}]{\sk{उपविशति} (upaviśati) — sits down}

\label{lesson:SA-C71-sits}



\begin{quote}
Before the new one: say the Sanskrit for he runs, then the Sanskrit for he walks.
\end{quote}

\subsection*{You'll want to know: \sk{उपविशति}}
\textbf{\sk{उपविशति}} (\emph{upaviśati}) — \textquotedblleft{}(he, she) sits down\textquotedblright{}, often on an \textbf{\sk{आसनम्}}.

\begin{cousinweb}
\emph{Upaviśati} is built from \emph{upa}, \textquotedblleft{}near\textquotedblright{}, and the root √viś, \textquotedblleft{}to enter, settle\textquotedblright{} (MW). It lives on as Hindi \emph{baiṭhnā}, \textquotedblleft{}to sit\textquotedblright{} (CDIAL 2245).
\end{cousinweb}

\begin{grammarlens}[title={What you've built}]
What a guest does first.
\end{grammarlens}

\subsection*{Your turn}
\begin{itemize}
\raggedright
  \item \emph{Say it:} \emph{upaviśati}
  \item \emph{Say it:} \emph{upaviśati}, once more
  \item \emph{Say it:} say \emph{atithiḥ upaviśati}
  \item \emph{Recall:} say the Sanskrit for he runs, then the Sanskrit for he walks, then say \emph{upaviśati} again
\end{itemize}

\begin{tcolorbox}[breakable,colback=teal!4,colframe=teal!35!black,title={Before you move on}]
What does \sk{उपविशति} mean? (\textquotedblleft{}Sits down\textquotedblright{}.) Say it once more.
\end{tcolorbox}
\section[\texorpdfstring{\sk{उत्तिष्ठति} (uttiṣṭhati)}{उत्तिष्ठति (uttiṣṭhati)}]{\sk{उत्तिष्ठति} (uttiṣṭhati) — stands up, gets up}

\label{lesson:SA-C71-standsup}



\begin{quote}
Before the new one: say the Sanskrit for he walks, then the Sanskrit for he sits down.
\end{quote}

\subsection*{You'll want to know: \sk{उत्तिष्ठति}}
\textbf{\sk{उत्तिष्ठति}} (\emph{uttiṣṭhati}) — \textquotedblleft{}(he, she) gets up\textquotedblright{}: \textbf{\sk{उत्}}, \textquotedblleft{}up\textquotedblright{}, with \textbf{\sk{तिष्ठति}}, \textquotedblleft{}stands\textquotedblright{}, which you know.

\begin{cousinweb}
\emph{Uttiṣṭhati} is built from \emph{ud}, \textquotedblleft{}up\textquotedblright{}, and \emph{tiṣṭhati}, \textquotedblleft{}stands\textquotedblright{}, from the root √sthā (CDIAL 1793).
\end{cousinweb}

\begin{grammarlens}[title={What you've built}]
Up plus stand.
\end{grammarlens}

\subsection*{Your turn}
\begin{itemize}
\raggedright
  \item \emph{Say it:} \emph{uttiṣṭhati}
  \item \emph{Say it:} \emph{uttiṣṭhati}, once more
  \item \emph{Say it:} say \emph{uttiṣṭhati}
  \item \emph{Recall:} say the Sanskrit for he walks, then the Sanskrit for he sits down, then say \emph{uttiṣṭhati} again
\end{itemize}

\begin{tcolorbox}[breakable,colback=teal!4,colframe=teal!35!black,title={Before you move on}]
What does \sk{उत्तिष्ठति} mean? (\textquotedblleft{}Stands up, gets up\textquotedblright{}.) Say it once more.
\end{tcolorbox}
\section[\texorpdfstring{\sk{पतति} (patati)}{पतति (patati)}]{\sk{पतति} (patati) — falls}

\label{lesson:SA-C71-falls}



\begin{quote}
Before the new one: say the Sanskrit for he sits down, then the Sanskrit for he stands up.
\end{quote}

\subsection*{You'll want to know: \sk{पतति}}
\textbf{\sk{पतति}} (\emph{patati}) — \textquotedblleft{}(he, she, it) falls\textquotedblright{}. A leaf, a \textbf{\sk{पत्रम्}}, \textbf{\sk{पतति}}.

\begin{cousinweb}
\emph{Patati} is the present of the root √pat, which first meant \textquotedblleft{}to fly\textquotedblright{} and then \textquotedblleft{}to fall\textquotedblright{} (CDIAL 7722). It is inherited: MW compares Greek \emph{petomai}, \textquotedblleft{}I fly\textquotedblright{}. It lives on as Hindi \emph{paṛnā}, \textquotedblleft{}to fall\textquotedblright{}.
\end{cousinweb}

\begin{grammarlens}[title={What you've built}]
What leaves and rain do.
\end{grammarlens}

\subsection*{Your turn}
\begin{itemize}
\raggedright
  \item \emph{Say it:} \emph{patati}
  \item \emph{Say it:} \emph{patati}, once more
  \item \emph{Say it:} say \emph{patraṁ patati}
  \item \emph{Recall:} say the Sanskrit for he sits down, then the Sanskrit for he stands up, then say \emph{patati} again
\end{itemize}

\begin{tcolorbox}[breakable,colback=teal!4,colframe=teal!35!black,title={Before you move on}]
What does \sk{पतति} mean? (\textquotedblleft{}Falls\textquotedblright{}.) Say it once more.
\end{tcolorbox}
